\documentclass[12pt,a4paper]{report}
\usepackage[utf8]{inputenc}
\usepackage[vietnamese]{babel}
\usepackage{geometry}
\geometry{left=3cm,right=2cm,top=2.5cm,bottom=2.5cm}

\begin{document}

    \chapter{GIỚI THIỆU}
    \section{Đặt vấn đề}
    \begin{itemize}
        \item Sự phát triển của IoT và nhu cầu về bảo mật thông minh trong kỷ nguyên số.
        \item Hạn chế của các phương pháp bảo mật truyền thống (khóa cơ, mật mã, thẻ từ).
        \item Xu hướng ứng dụng sinh trắc học và Edge AI vào hệ thống nhúng.
    \end{itemize}

    \section{Mục tiêu của đề tài}
    \begin{itemize}
        \item \textbf{Mục tiêu tổng quát:} Xây dựng hệ thống khóa thông minh nhận diện khuôn mặt hoạt động ổn định trên phần cứng chi phí thấp.
        \item \textbf{Mục tiêu cụ thể:} Thiết kế phần cứng (ESP32-S3 Cam), triển khai mô hình nhận diện gọn nhẹ, xây dựng ứng dụng quản lý (Mobile/Web), tối ưu hóa hiệu năng.
    \end{itemize}

    \section{Đối tượng và phạm vi nghiên cứu}
    \begin{itemize}
        \item \textbf{Đối tượng:} Board ESP32-S3 Cam, TinyML, cơ chế điều khiển khóa.
        \item \textbf{Phạm vi:} Nhận diện trong môi trường thực tế (30-80cm, ánh sáng văn phòng/gia đình), quản lý người dùng cục bộ hoặc Server mini.
    \end{itemize}

    \section{Ý nghĩa khoa học và thực tiễn}
    \begin{itemize}
        \item \textbf{Khoa học:} Thử nghiệm khả năng tính toán của ESP32-S3 với các tác vụ thị giác máy tính.
        \item \textbf{Thực tiễn:} Giải pháp khóa cửa thông minh chi phí tối ưu cho phòng trọ, văn phòng nhỏ.
    \end{itemize}

    \section{Cấu trúc của báo cáo}

    \chapter{CƠ SỞ LÝ THUYẾT VÀ CÔNG NGHỆ SỬ DỤNG}
    \section{Cơ sở lý thuyết về thị giác máy tính và Nhận diện khuôn mặt}
    \subsection{Bài toán phát hiện khuôn mặt (Face Detection)}
    \subsection{Bài toán trích xuất đặc trưng và nhận diện (Face Recognition)}
    \subsection{Tổng quan về TinyML}

    \section{Công nghệ phần cứng sử dụng}
    \subsection{ESP32-S3 CAM}
    \subsection{Các linh kiện ngoại vi (Relay, Buck converter, cơ cấu khóa)}

    \section{Công nghệ phần mềm và Thư viện}
    \subsection{ESP-IDF / Arduino IDE}
    \subsection{ESP-WHO / ESP-DL}
    \subsection{Nền tảng quản lý (Flutter, NestJS/Prisma)}

    \chapter{PHÂN TÍCH VÀ THIẾT KẾ HỆ THỐNG}
    \section{Phân tích yêu cầu hệ thống}
    \subsection{Yêu cầu chức năng}
    \subsection{Yêu cầu phi chức năng}

    \section{Sơ đồ kiến trúc hệ thống}
    \section{Các biểu đồ UML}
    \subsection{Biểu đồ Use Case}
    \subsection{Biểu đồ Hoạt động (Activity Diagram)}
    \subsection{Biểu đồ Trình tự (Sequence Diagram)}

    \section{Thiết kế cơ sở dữ liệu}
    \subsection{Lưu trữ trên Flash (Cục bộ ESP32)}
    \subsection{Lưu trữ trên Server (ERD)}

    \chapter{TRIỂN KHAI THỬ NGHIỆM VÀ ĐÁNH GIÁ}
    \section{Cấu hình môi trường và triển khai}
    \section{Kịch bản kiểm thử (Testcases)}
    \section{Kết quả triển khai và Phân tích đánh giá}
    \subsection{Kết quả đạt được}
    \subsection{Bảng thống kê hiệu năng (FAR, FRR, Tốc độ xử lý)}
    \subsection{Phân tích hạn chế}

    \chapter{KẾT LUẬN}
    \section{Các đóng góp chính của đề tài}
    \section{Hướng phát triển trong tương lai}

    \chapter*{TÀI LIỆU THAM KHẢO}

\end{document}